\documentclass[11pt]{article}
\usepackage[a4paper,margin=2.4cm]{geometry}
\usepackage{amsmath,amssymb,amsthm}
\usepackage{mathtools}
\usepackage{enumitem}
\usepackage[colorlinks=true,linkcolor=blue,urlcolor=blue]{hyperref}
\newcommand{\R}{\mathbb{R}}
\newcommand{\E}{\mathbb{E}}
\newcommand{\dd}{\,\mathrm{d}}
\newtheorem{remark}{Remark}

\newcommand{\impl}{\mathrm{impl}}
\newcommand{\mkt}{\mathrm{mkt}}
\newcommand{\BS}{\mathrm{BS}}
\newcommand{\Vega}{\mathcal{V}}

\title{\textbf{Lecture 3 --- Implied Volatility}\\[2pt]
\large Study notes: definition, existence and uniqueness, root-finding (bisection, Newton),
the volatility surface, and the deficiencies of Black--Scholes}
\author{Computational Finance (WI4151), TU Delft --- F.\ Fang}
\date{}

\begin{document}
\maketitle
\vspace{-1.2em}

\begin{abstract}
\noindent These notes work through Lecture 3 in full. Implied volatility is the value of the
single free Black--Scholes parameter $\sigma$ that reproduces an observed market option price.
The lecture has two halves. First, the \emph{theory}: $\sigma\mapsto C^{\BS}(\sigma)$ is
continuous, strictly increasing (because $\Vega>0$) with explicit limits as $\sigma\to0^+$ and
$\sigma\to\infty$, so the inverse exists and is unique on the admissible price interval. Second,
the \emph{numerics}: bisection (robust, linear) and Newton (fragile but quadratic), with the
key practical result that starting Newton at the inflection point $\hat\sigma$ of $C(\sigma)$
gives \emph{global} monotone convergence. Finally, the empirical observation that implied
volatility is not flat across strike and maturity --- the smile/skew/surface --- which is the
diagnostic that constant-volatility Black--Scholes is misspecified, and the motivation for the
local-, stochastic-volatility and jump models of later lectures. Everything is derived; where I
fill gaps not on the slides (notably the closed form of $\Vega$, the second-derivative identity,
and the put treatment) I say so explicitly.
\end{abstract}

\tableofcontents

\section{Setup: the four knowns and the one unknown}

We work inside the Black--Scholes model. Under the risk-neutral measure $\mathbb{Q}$ the stock
solves $\dd S_t=rS_t\dd t+\sigma S_t\dd W_t^{\mathbb{Q}}$, and the time-$t$ price of a European
call with strike $K$ and maturity $T$ is the closed form
\begin{equation}\label{eq:bs}
C^{\BS}(S_t,t;\sigma)=S_t\,N(d_1)-Ke^{-r\tau}N(d_2),\qquad \tau:=T-t,
\end{equation}
with
\begin{equation}\label{eq:d1d2}
d_1=\frac{\ln(S_t/K)+(r+\tfrac12\sigma^2)\tau}{\sigma\sqrt{\tau}},\qquad
d_2=\frac{\ln(S_t/K)+(r-\tfrac12\sigma^2)\tau}{\sigma\sqrt{\tau}}=d_1-\sigma\sqrt{\tau},
\end{equation}
and $N$ the standard normal CDF, $N'(x)=\tfrac{1}{\sqrt{2\pi}}e^{-x^2/2}$.

The formula has five inputs $(\tau,K,r,S_t,\sigma)$. Four of them are directly observable:
the time to maturity $\tau=T-t$, the strike $K$, the risk-free rate $r$, and the current spot
$S_t$. The fifth, the volatility $\sigma$, is \emph{not} observable: it is a property of the
unobserved future, not a quoted number. Black--Scholes treats it as a constant model parameter,
but the market does not hand it to us.

\paragraph{The inverse problem.} The market \emph{does} hand us option prices $C^{\mkt}$. So we
turn \eqref{eq:bs} around. Holding $(\tau,K,r,S_t)$ fixed at their observed values, regard the
call price as a scalar function of the one free parameter,
\[
\sigma\;\longmapsto\;C^{\BS}(\sigma),
\]
and ask: which $\sigma$ reproduces the quote? The \emph{implied volatility} $\sigma^{\impl}$ is
defined as the solution of
\begin{equation}\label{eq:def-iv}
V^{\BS}(\sigma^{\impl};r,\tau,K,S_t)-V^{\mkt}=0,
\end{equation}
i.e. the volatility ``implied'' by the market price through the Black--Scholes lens. It is the
market's quote expressed in the units of the model. We treat the call below; the put is identical
after one remark (Section~\ref{sec:put}).

\begin{remark}
Implied volatility is \emph{not} a forecast of realised volatility and \emph{not} a directly
measured quantity. It is a change of variable: a one-to-one re-quoting of a price as a vol,
using a model we already suspect is wrong. Its usefulness is precisely that, were Black--Scholes
true, $\sigma^{\impl}$ would be the same constant for every $(K,T)$. The extent to which it is
\emph{not} constant is a direct readout of the model's misspecification.
\end{remark}

\section{Existence and uniqueness of the implied volatility}

For \eqref{eq:def-iv} to define $\sigma^{\impl}$ unambiguously we need: (i) $C^{\BS}(\sigma)$
strictly monotone in $\sigma$ (so the inverse is single-valued), and (ii) the market price to lie
in the range of $C^{\BS}$. Both follow from elementary properties of \eqref{eq:bs}.

\subsection{Vega is strictly positive}

The \emph{vega} is the sensitivity of the price to volatility,
\[
\Vega:=\frac{\partial C}{\partial\sigma}.
\]
Differentiate \eqref{eq:bs}. By the chain rule, treating $d_1,d_2$ as functions of $\sigma$,
\begin{equation}\label{eq:vega-chain}
\frac{\partial C}{\partial\sigma}
= S_t N'(d_1)\,\frac{\partial d_1}{\partial\sigma}
 - Ke^{-r\tau}N'(d_2)\,\frac{\partial d_2}{\partial\sigma}.
\end{equation}
From $d_2=d_1-\sigma\sqrt{\tau}$ we get $\dfrac{\partial d_2}{\partial\sigma}
=\dfrac{\partial d_1}{\partial\sigma}-\sqrt{\tau}$. The key simplification is the identity
\begin{equation}\label{eq:vega-identity}
S_t\,N'(d_1)=Ke^{-r\tau}\,N'(d_2),
\end{equation}
which I prove below. Substituting \eqref{eq:vega-identity} into \eqref{eq:vega-chain}, the two
terms carrying $\partial d_1/\partial\sigma$ have equal coefficients $S_tN'(d_1)=Ke^{-r\tau}N'(d_2)$
and cancel:
\[
\frac{\partial C}{\partial\sigma}
= S_tN'(d_1)\frac{\partial d_1}{\partial\sigma}
 - Ke^{-r\tau}N'(d_2)\Big(\frac{\partial d_1}{\partial\sigma}-\sqrt{\tau}\Big)
= Ke^{-r\tau}N'(d_2)\,\sqrt{\tau}.
\]
Using \eqref{eq:vega-identity} once more to rewrite the prefactor in terms of $S_t$ and $d_1$,
\begin{equation}\label{eq:vega}
\boxed{\;\Vega=\frac{\partial C}{\partial\sigma}=S_t\,N'(d_1)\,\sqrt{\tau}
=\frac{S_t\sqrt{\tau}}{\sqrt{2\pi}}\,e^{-\frac12 d_1^2}\;>0\;}
\qquad(\sigma>0,\ \tau>0).
\end{equation}
Every factor on the right is strictly positive, so $\Vega>0$. Hence $C^{\BS}(\sigma)$ is
\emph{strictly increasing} in $\sigma$ on $(0,\infty)$. This is the structural fact that makes the
implied volatility well posed: a strictly monotone continuous function has a continuous inverse,
so \emph{at most one} $\sigma^{\impl}$ can solve \eqref{eq:def-iv}.

\paragraph{Proof of identity \eqref{eq:vega-identity}.} (Not on the slides; supplied here because
Fang is exactly the kind of examiner who asks ``why does that identity hold?'') Write the ratio
\[
\frac{N'(d_1)}{N'(d_2)}=\exp\!\Big(-\tfrac12(d_1^2-d_2^2)\Big)
=\exp\!\Big(-\tfrac12(d_1-d_2)(d_1+d_2)\Big).
\]
Now $d_1-d_2=\sigma\sqrt{\tau}$ and $d_1+d_2=\dfrac{2\ln(S_t/K)+2r\tau}{\sigma\sqrt{\tau}}$
(the $\pm\tfrac12\sigma^2\tau$ terms cancel in the sum). Therefore
\[
-\tfrac12(d_1-d_2)(d_1+d_2)
=-\tfrac12\,\sigma\sqrt{\tau}\cdot\frac{2\ln(S_t/K)+2r\tau}{\sigma\sqrt{\tau}}
=-\ln(S_t/K)-r\tau,
\]
so $\dfrac{N'(d_1)}{N'(d_2)}=e^{-\ln(S_t/K)-r\tau}=\dfrac{K}{S_t}\,e^{-r\tau}$, i.e.
$S_tN'(d_1)=Ke^{-r\tau}N'(d_2)$. \hfill$\square$

\begin{remark}[Vega's shape]
$\Vega=S_t\sqrt{\tau}\,N'(d_1)$ is a Gaussian bump in $d_1$, maximised when $d_1=0$, i.e. near the
money. It decays to zero for deep in- or out-of-the-money options ($|d_1|\to\infty$). This will
matter twice below: it is why Newton can stall in the wings (tiny denominator), and it is the
reason the implied-vol \emph{inversion} is ill-conditioned far from the money --- a small price
error is divided by a near-zero vega and blows up the implied vol.
\end{remark}

\subsection{The admissible price interval}

Monotonicity gives \emph{at most} one solution. Existence requires the quote to lie in the range
of $C^{\BS}$, which we now pin down by taking the two limits in $\sigma$.

\paragraph{Lower limit, $\sigma\to0^+$.} As $\sigma\to0^+$ the sign of $d_1,d_2$ is governed by
the sign of $\ln(S_t/K)+r\tau$ (the volatility-independent part of the numerator), since
$\sigma\sqrt{\tau}\to0$ in the denominator amplifies that sign:
\begin{itemize}[nosep]
\item If $\ln(S_t/K)+r\tau>0$: $d_1,d_2\to+\infty$, $N(d_1),N(d_2)\to1$, so
      $C(\sigma)\to S_t-Ke^{-r\tau}>0$.
\item If $\ln(S_t/K)+r\tau<0$: $d_1,d_2\to-\infty$, $N(d_1),N(d_2)\to0$, so $C(\sigma)\to0$.
\item If $\ln(S_t/K)+r\tau=0$: $d_1,d_2\to0$, $N\to\tfrac12$, so
      $C(\sigma)\to\tfrac12(S_t-Ke^{-r\tau})=0$, the last equality because the case condition is
      exactly $S_t=Ke^{-r\tau}$.
\end{itemize}
The three branches collapse to one statement,
\begin{equation}\label{eq:lower}
\lim_{\sigma\to0^+}C(\sigma)=\max\!\big(S_t-Ke^{-r\tau},\,0\big),
\end{equation}
which is the no-arbitrage \emph{intrinsic value} of the call (its forward-discounted moneyness,
floored at zero). This is the natural floor: a zero-vol stock grows deterministically at $r$, and
the call is worth its discounted in-the-money amount or nothing.

\paragraph{Upper limit, $\sigma\to\infty$.} As $\sigma\to\infty$, the $+\tfrac12\sigma^2\tau$ term
sends $d_1\to+\infty$ while the $-\tfrac12\sigma^2\tau$ term sends $d_2\to-\infty$. Hence
$N(d_1)\to1$, $N(d_2)\to0$, and
\begin{equation}\label{eq:upper}
\lim_{\sigma\to\infty}C(\sigma)=S_t.
\end{equation}
Economically: with infinite volatility the downside is unchanged (payoff floored at $0$) but the
upside is unbounded, and the option's value rises to the whole spot $S_t$ (the discounted strike
becomes negligible against the dispersion). It is a supremum, not attained.

\paragraph{Conclusion.} Combining \eqref{eq:vega}, \eqref{eq:lower}, \eqref{eq:upper}: for every
$\sigma\ge0$,
\begin{equation}\label{eq:bounds}
\max\!\big(S_t-Ke^{-r\tau},0\big)\;\le\;C(\sigma)\;<\;S_t,
\end{equation}
and $C$ is a continuous strictly increasing bijection from $[0,\infty)$ onto
$\big[\max(S_t-Ke^{-r\tau},0),\,S_t\big)$. Therefore:
\begin{quote}
\emph{If} $C^{\mkt}\in\big[\max(S_t-Ke^{-r\tau},0),\,S_t\big)$, then there is a \emph{unique}
$\sigma^{\impl}\ge0$ with $C(\sigma^{\impl})=C^{\mkt}$.
\end{quote}
A quote outside this band violates a static no-arbitrage bound (a call worth less than intrinsic,
or more than the stock) and admits no real implied volatility --- which in practice flags a stale
or erroneous quote, or a quote at the bid/ask edge.

\section{Numerical inversion}

The equation $F(\sigma):=C(\sigma)-C^{\mkt}=0$ has no closed-form solution (it sits behind the
nonlinear $N(d_1),N(d_2)$), so we root-find numerically. Two methods are on the slides.

\subsection{Bisection}

\paragraph{Idea.} If $F$ is continuous and changes sign on $[\sigma_a,\sigma_b]$, i.e.
$F(\sigma_a)F(\sigma_b)<0$, the intermediate value theorem guarantees a root inside. Evaluate the
midpoint $\sigma_m=\tfrac12(\sigma_a+\sigma_b)$; the sign of $F(\sigma_m)$ agrees with exactly one
endpoint, so one of $[\sigma_a,\sigma_m]$, $[\sigma_m,\sigma_b]$ still brackets the root. Replace
the bracket by that half and repeat.

\paragraph{Pseudocode (as on the slide).}
\begin{verbatim}
N <- 1
while N <= NMAX:
    c <- (a+b)/2                      # midpoint
    if F(c)=0 or (b-a)/2 < TOL:       # converged
        output c; stop
    N <- N+1
    if sign(F(c)) = sign(F(a)): a <- c    # root in [c,b]
    else:                       b <- c    # root in [a,c]
\end{verbatim}

\paragraph{Why it always works here.} By \eqref{eq:bounds} and monotonicity, an initial bracket
exists whenever the quote is admissible: take $\sigma_a$ tiny (then $F<0$ once
$C^{\mkt}>$ intrinsic) and $\sigma_b$ large (then $F>0$ once $C^{\mkt}<S_t$). Monotonicity of $C$
guarantees a single sign change, so bisection cannot be confused by spurious roots.

\paragraph{Convergence.} The bracket halves each step: after $n$ iterations the root is localised
to width $(\sigma_b-\sigma_a)/2^{\,n}$. So the error obeys
$|\sigma_n-\sigma^*|\le 2^{-n}(\sigma_b-\sigma_a)$ --- \emph{linear} convergence with rate $1/2$,
one bit of accuracy per iteration. To reach tolerance $\varepsilon$ needs about
$\log_2\!\big((\sigma_b-\sigma_a)/\varepsilon\big)$ steps.

\paragraph{Trade-off.} Advantage: needs no derivative and no good initial guess --- only a sign-
changing bracket, which we can always produce. It is globally convergent and bullet-proof.
Disadvantage: only linear, so it is slow when high accuracy is wanted.

\subsection{Newton--Raphson}

\paragraph{Derivation.} To solve $F(\sigma^*)=0$, Taylor-expand $F$ about the current iterate
$\sigma_n$ for a small step $\delta$:
\[
F(\sigma_n+\delta)=F(\sigma_n)+\delta F'(\sigma_n)+\mathcal O(\delta^2).
\]
Drop the quadratic remainder and demand the linear model vanish, $F(\sigma_n)+\delta F'(\sigma_n)=0$.
Solving for $\delta$ and setting $\sigma_{n+1}=\sigma_n+\delta$ gives the iteration
\begin{equation}\label{eq:newton}
\sigma_{n+1}=\sigma_n-\frac{F(\sigma_n)}{F'(\sigma_n)}
=\sigma_n-\frac{C(\sigma_n)-C^{\mkt}}{\Vega(\sigma_n)}.
\end{equation}
Here $F'=C'=\Vega$ from \eqref{eq:vega}, so each step needs one price and one vega evaluation ---
both cheap, and both already available in closed form. Geometrically, $\sigma_{n+1}$ is where the
tangent to $F$ at $\sigma_n$ crosses zero.

\paragraph{Local convergence theorem.} If $F\in C^2$ and there is a simple root $\sigma^*$ with
$F(\sigma^*)=0$, $F'(\sigma^*)\neq0$, then there is $\delta>0$ such that for any start with
$|\sigma_0-\sigma^*|<\delta$ the iterates are well defined, $\sigma_n\to\sigma^*$, and there is a
constant $C$ with
\begin{equation}\label{eq:quad}
|\sigma_{n+1}-\sigma^*|\le C\,|\sigma_n-\sigma^*|^2.
\end{equation}
This is \emph{quadratic} convergence: the number of correct digits roughly doubles each step. For
the IV problem $F'=\Vega>0$ at any interior root, so the simple-root hypothesis holds.

\paragraph{Why quadratic.} Subtract $\sigma^*$ from \eqref{eq:newton} and use $F(\sigma^*)=0$:
\begin{align*}
\sigma_{n+1}-\sigma^*
&=\sigma_n-\sigma^*-\frac{F(\sigma_n)}{F'(\sigma_n)}
=\sigma_n-\sigma^*-\frac{F(\sigma_n)-F(\sigma^*)}{F'(\sigma_n)}\\
&=\sigma_n-\sigma^*-\frac{(\sigma_n-\sigma^*)F'(\sigma_n)+\mathcal O\big((\sigma_n-\sigma^*)^2\big)}
{F'(\sigma_n)}
=\mathcal O\big((\sigma_n-\sigma^*)^2\big),
\end{align*}
where the third equality is the Taylor expansion
$F(\sigma_n)-F(\sigma^*)=(\sigma_n-\sigma^*)F'(\sigma_n)-\tfrac12(\sigma_n-\sigma^*)^2F''(\xi)+\cdots$
rearranged so the linear part cancels $(\sigma_n-\sigma^*)$. The leading surviving term is
quadratic, with constant $\approx \tfrac12|F''/F'|$ near $\sigma^*$.

\paragraph{Trade-off.} Advantage: much faster than bisection (quadratic vs.\ linear). Disadvantage:
\emph{local} --- a poor $\sigma_0$ can overshoot, diverge, or step to negative volatility, and in
the deep wings $\Vega(\sigma_n)\approx0$ makes the update \eqref{eq:newton} explode. So Newton is
fast where it works but is not robust by itself.

\subsection{Globalising Newton: start at the inflection point $\hat\sigma$}

The slides give a clean cure for Newton's fragility: a special starting point from which
convergence is \emph{global and monotone}. The construction rests on the convexity structure of
$C(\sigma)$.

\paragraph{Second derivative.} Differentiating $\Vega=S_t\sqrt{\tau}\,N'(d_1)$ in $\sigma$ and
simplifying (using $\partial d_1/\partial\sigma=-d_2/\sigma$, a standard reduction) gives the
slide's formula
\begin{equation}\label{eq:d2C}
\frac{\partial^2 C}{\partial\sigma^2}
=\frac{S_t\sqrt{\tau}}{\sqrt{2\pi}}\,e^{-\frac12 d_1^2}\cdot\frac{d_1 d_2}{\sigma},
\qquad\sigma>0.
\end{equation}
The prefactor is positive, so the sign of $\partial^2 C/\partial\sigma^2$ is the sign of $d_1d_2$.
Now
\[
d_1 d_2
=\frac{\ln(S_t/K)+(r+\tfrac12\sigma^2)\tau}{\sigma\sqrt\tau}\cdot
 \frac{\ln(S_t/K)+(r-\tfrac12\sigma^2)\tau}{\sigma\sqrt\tau}
=\frac{\big(\ln(S_t/K)+r\tau\big)^2-\big(\tfrac12\sigma^2\tau\big)^2}{\sigma^2\tau},
\]
a difference of squares. This vanishes precisely when
$\big(\ln(S_t/K)+r\tau\big)^2=\tfrac14\sigma^4\tau^2$, i.e. at
\begin{equation}\label{eq:sigmahat}
\sigma=\hat\sigma:=\sqrt{\frac{2\,\big|\ln(S_t/K)+r\tau\big|}{\tau}}.
\end{equation}
For $0\le\sigma<\hat\sigma$ the squared moneyness term dominates, $d_1d_2>0$, so
$\partial^2 C/\partial\sigma^2>0$ and $C$ is \emph{convex}. For $\sigma>\hat\sigma$ we have
$d_1d_2<0$, so $\partial^2 C/\partial\sigma^2<0$ and $C$ is \emph{concave}. Thus $\hat\sigma$ is
the unique inflection point of $C(\sigma)$, and --- crucially --- it is where $\Vega=C'$ attains
its \emph{maximum} (the slope is largest exactly where curvature changes sign).

\paragraph{Global monotone convergence from $\sigma_0=\hat\sigma$.} Suppose the true root satisfies
$\sigma^*>\sigma_0=\hat\sigma$ (the case $\sigma^*<\hat\sigma$ is symmetric). By the mean value
theorem applied to $F$ on $[\sigma_0,\sigma^*]$, there is $\xi_0\in(\sigma_0,\sigma^*)$ with
$F(\sigma_0)-F(\sigma^*)=(\sigma_0-\sigma^*)F'(\xi_0)$, hence from \eqref{eq:newton}
\[
\sigma_1-\sigma^*
=(\sigma_0-\sigma^*)-\frac{(\sigma_0-\sigma^*)F'(\xi_0)}{F'(\sigma_0)}
\;\Longrightarrow\;
\frac{\sigma_1-\sigma^*}{\sigma_0-\sigma^*}=1-\frac{F'(\xi_0)}{F'(\sigma_0)}.
\]
Because $\sigma_0=\hat\sigma$ \emph{maximises} $F'=\Vega$, we have $0<F'(\xi_0)\le F'(\sigma_0)$, so
the ratio lies in $(0,1)$. Therefore $\sigma_1-\sigma^*$ has the same sign as $\sigma_0-\sigma^*$
but smaller magnitude: $\sigma_0<\sigma_1<\sigma^*$. At the next step $\sigma_1$ lies on the
concave branch where $F'=\Vega$ is positive and \emph{decreasing}, and the same MVT argument with
$\sigma_1<\xi_1<\sigma^*$ gives the ratio again in $(0,1)$, so $\sigma_1<\sigma_2<\sigma^*$.
Inductively $\sigma_n<\sigma_{n+1}<\sigma^*$: the iterates form an increasing sequence bounded
above by $\sigma^*$, hence converge, and the limit is the root. The error decreases monotonically
to zero and, by \eqref{eq:quad}, quadratically near the end.

\begin{remark}
The starting point \eqref{eq:sigmahat} is the trick that makes Newton both fast and safe for this
specific problem: the maximal-slope property of $\hat\sigma$ is exactly what forces the
contraction ratio $1-F'(\xi)/F'(\sigma_0)$ into $(0,1)$ on the first step, after which the
iterates stay on the ``good'' (concave, monotone) side. This is a problem-specific globalisation,
not a general Newton guarantee. In practice one still caps iterations and falls back to bisection
if a step leaves the admissible interval.
\end{remark}

\subsection{Method comparison}

\begin{center}
\begin{tabular}{lll}
\hline
 & \textbf{Bisection} & \textbf{Newton (from $\hat\sigma$)}\\
\hline
Order of convergence & linear, rate $\tfrac12$ & quadratic\\
Needs derivative? & no & yes ($\Vega$, closed form)\\
Needs good start? & no, only a bracket & yes; $\hat\sigma$ supplies it\\
Robustness & global, fail-safe & global here by the $\hat\sigma$ argument\\
Cost per step & one price eval & one price + one vega eval\\
\hline
\end{tabular}
\end{center}

A common production choice is a hybrid: a few bisection steps to localise, then switch to Newton
for the final digits; or Newton from $\hat\sigma$ with a bisection safety net.

\section{The put, and a sanity check via parity}\label{sec:put}

The put is handled identically. By put--call parity $C-P=S_t-Ke^{-r\tau}$, which is
$\sigma$-independent, the put price is $P(\sigma)=C(\sigma)-(S_t-Ke^{-r\tau})$. Differentiating,
\[
\frac{\partial P}{\partial\sigma}=\frac{\partial C}{\partial\sigma}=\Vega=S_t\sqrt\tau\,N'(d_1)>0,
\]
so the put has the \emph{same} positive vega and the same strict monotonicity; its implied
volatility is unique on the analogous admissible band, and an arbitrage-free call and put at the
same $(K,T)$ return the \emph{identical} implied volatility. This is why one may quote one IV per
$(K,T)$ regardless of call/put. The same bisection/Newton machinery applies verbatim with $C$
replaced by $P$.

\section{The implied volatility surface}

If Black--Scholes held exactly, $\sigma^{\impl}$ recovered from \eqref{eq:def-iv} would be a single
constant, identical across every strike $K$ and maturity $T$ --- because $\sigma$ is, in the model,
one fixed number. The empirical fact is that it is not. Inverting market quotes produces a function
\[
(K,T)\longmapsto\sigma^{\impl}(K,T),
\]
the \emph{implied volatility surface}. Its non-flatness is the central diagnostic of this lecture.

\paragraph{Empirical example (FTSE 100, 22 Aug 2001).} With $S_0=5420.3$, $r=0.05$,
$T=\tfrac13$\,yr, inverting the quoted call prices across strikes $5125,\dots,5825$ yields implied
volatilities that \emph{fall} monotonically with strike (from $\approx0.190$ at $K=5125$ to
$\approx0.174$ at $K=5825$) instead of staying flat. Under Black--Scholes they should coincide;
they do not, so the market does not obey the model's constant-vol assumption.

\subsection{Shapes in the strike direction}

Holding $T$ fixed, $K\mapsto\sigma^{\impl}(K,T)$ takes characteristic shapes:
\begin{itemize}[nosep]
\item \textbf{Smile.} A convex U: implied vol is lowest near the money and higher for both deep
      ITM and deep OTM strikes. Reading it through the model, the market prices \emph{both} tails
      richer than a lognormal --- consistent with a return distribution with heavier tails
      (excess kurtosis) than the Gaussian Black--Scholes assumes.
\item \textbf{Skew (smirk).} A monotone decreasing curve: implied vol is high for low strikes
      (downside puts) and decays for high strikes. This is the typical equity-index shape and
      encodes a left-skewed (negatively skewed) risk-neutral distribution: crash risk is priced
      richer, so out-of-the-money puts carry higher implied vol.
\item \textbf{Hockey stick.} A skew on the left blending into a smile on the right --- a
      combination of the two, with a rich left tail and a modestly rich right tail.
\end{itemize}

\subsection{Term structure in the maturity direction}

Holding $K$ fixed (say at the money), $T\mapsto\sigma^{\impl}(K,T)$ is the \emph{term structure}.
The slide contrasts two models for $\dd S/S=\mu\dd t+\sigma(t)\dd W_t$:
\begin{itemize}[nosep]
\item Constant volatility $\sigma_*$: the ATM implied term structure is flat in $T$.
\item Time-dependent $\sigma(t)$: the implied vol at maturity $T$ is the root-mean-square of the
      instantaneous vol, $\sigma^{\impl}(T)^2=\tfrac1T\int_0^T\sigma(t)^2\dd t$, so a rising
      $\sigma(t)$ produces an upward-sloping implied term structure. (This is the cleanest way to
      see that the BS implied vol over $[0,T]$ aggregates the path of instantaneous variance into
      a single average.)
\end{itemize}

\subsection{The full surface}

Combining both directions gives $\sigma^{\impl}(K,T)$ over a 2-D domain. The slides show the
generic pattern: a pronounced \emph{smile at short maturities} (the wings are most curved when
there is little time for the central limit theorem to Gaussianise the return), flattening into a
milder \emph{skew at longer maturities}. A flat plane would mean Black--Scholes; the curvature and
tilt of the actual surface is the market's correction to it, and any candidate model is judged by
how well it can reproduce this surface (calibration).

\section{Why constant-vol Black--Scholes fails}

The surface is a symptom; the lecture closes by naming the broken assumptions and the families of
models that repair them.

\paragraph{Assumptions that fail in reality.}
\begin{itemize}[nosep]
\item \textbf{Continuous trading.} Delta hedging is assumed continuous-time and frictionless; real
      rebalancing is discrete, leaving hedging error.
\item \textbf{No transaction costs.} Real hedging incurs costs, which feed back into prices.
\item \textbf{Gaussian returns.} BS log-returns are normal; empirical returns are heavy-tailed and
      negatively skewed --- exactly what the smile/skew reports.
\item \textbf{Constant volatility.} Realised volatility is not constant: it clusters (calm and
      turbulent regimes persist), so a single $\sigma$ cannot fit prices across $(K,T)$
      simultaneously.
\end{itemize}

\paragraph{Repair families (introduced here, developed later).}
\begin{itemize}[nosep]
\item \textbf{Term-structure volatility:} $\dd S/S=\mu\dd t+\sigma(t)\dd W_t$ --- deterministic
      $\sigma(t)$ fits the maturity dimension but, being deterministic, still gives a flat smile
      in strike at each $T$.
\item \textbf{Local volatility:} $\dd S/S=\mu\dd t+\sigma(S,t)\dd W_t$ --- volatility a
      deterministic function of spot and time, calibrated (Dupire) to reproduce the \emph{entire}
      surface exactly, at the cost of unrealistic forward smile dynamics.
\item \textbf{Stochastic volatility:} $\dd S/S=\mu\dd t+\sigma(W^v,t)\dd W^S_t$ with volatility
      itself driven by a second Brownian motion --- e.g. Heston (Lecture 7). A genuinely random
      variance generates smiles/skews endogenously and is the workhorse for the COS method later.
\item \textbf{Jump models:} adding a jump component (Merton, Lecture 8) injects discontinuities
      that fatten the short-maturity tails, fixing the steep short-dated smile that diffusion
      alone cannot.
\end{itemize}
These are precisely the models the rest of the course prices, and the implied-vol surface is the
target they must hit.

\section{Distinguish: implied vs.\ local vs.\ stochastic volatility}

A point Fang may probe, since the slide lists $\sigma(t),\sigma(S,t),\sigma(W^v,t)$ side by side.
\begin{itemize}[nosep]
\item \textbf{Implied volatility} is a \emph{model-inversion output}, not a model input: a single
      number per quote, defined only through the Black--Scholes formula. It is the market price
      re-expressed in vol units. It says nothing on its own about the true dynamics.
\item \textbf{Local volatility} $\sigma(S,t)$ is a \emph{model input}: a deterministic surface
      that, by Dupire, can be chosen to match all implied vols at once. It is the
      ``effective'' instantaneous vol consistent with today's surface.
\item \textbf{Stochastic volatility} $\sigma(W^v,t)$ is also a model input but \emph{random},
      with its own driving noise (and, in Heston, correlation $\rho$ with the spot's noise). It
      explains \emph{why} the surface has the shape and dynamics it does, rather than merely
      fitting it.
\end{itemize}
In one line: implied vol describes the data, local vol fits it deterministically, stochastic vol
explains it.

\section*{Missing / not on the slides (flagged honestly)}
\addcontentsline{toc}{section}{Missing / not on the slides}
\begin{itemize}[nosep]
\item The closed form $\Vega=S_t\sqrt\tau\,N'(d_1)$ \eqref{eq:vega} is \emph{derived} here; the
      slides assert $\Vega>0$ and the cancellation but stop short of the boxed formula. The proof
      of the identity \eqref{eq:vega-identity} is mine; the slide only states it ``follows from
      $d_2=d_1-\sigma\sqrt\tau$ and the exponential form of $N'$''.
\item The reduction $\partial d_1/\partial\sigma=-d_2/\sigma$ used to get \eqref{eq:d2C} is stated
      without derivation on the slides; I quote it.
\item The put treatment via parity (Section~\ref{sec:put}) is my addition; the slides only say ``an
      analogous treatment applies''.
\item The RMS reading of the term structure,
      $\sigma^{\impl}(T)^2=\tfrac1T\int_0^T\sigma(t)^2\dd t$, is the standard interpretation of the
      $\sigma(t)$ figure but is not written on the slide.
\item The Dupire/local-vol and Heston cross-references are forward pointers to Lectures 7--8, not
      content of Lecture 3.
\end{itemize}

\end{document}
